\documentclass{RequiredFiles/NorthumbriaPhdThesis}
\usepackage[utf8]{inputenc}
\usepackage[T1]{fontenc}
\usepackage{geometry}
\usepackage{graphicx}
\usepackage{float}
\usepackage{hyperref}
\usepackage[english,indonesian]{babel}

\title{Optimasi Pemrosesan Template LaTeX Menggunakan Pendekatan Berbasis AI}
\author{Academic Support Agent}
\date{\today}

\begin{document}
\maketitle

\begin{abstract}

\end{abstract}

\chapter{PENDAHULUAN}


\chapter{TINJAUAN PUSTAKA}


% Original content (commented out):
% \documentclass[twoside, openright]{RequiredFiles/NorthumbriaPhdThesis}

%%%
% Meta Data (Used on Title Pages)
%%%
\title{Northumbria University \LaTeX~Thesis Template}
\author{Full Name}
\faculty{Faculty Name}
\wordcount{83,564}
\degreetype{Doctor of Philosophy}
\degreetitle{PhD}

\begin{document}

%%%
% Begin Front Matter
%%%
\frontmatter
\maketitle
\input{FrontMatter/FrontMatter.tex}

%%%
% Begin Main Matter
%%%
\mainmatter
\pagestyle{fancy}
\input{MainMatter/Chapters.tex}

%%%
% Begin Back Matter
%%%
\input{BackMatter/BackMatter.tex}

\end{document}

\end{document}